\documentclass[10pt,a4paper]{amsart}
\usepackage[margin=19mm]{geometry}
\usepackage{amsmath,amssymb,mathtools}
\usepackage[scr=rsfs,cal=euler]{mathalfa}
\usepackage{xfrac}
\usepackage[unicode,hidelinks,bookmarks=false]{hyperref}
\newcommand{\ie}{i.e.\ }
\newcommand{\cf}{cf.\ }
\newcommand{\resp}{respectively\ }
\newcommand{\Subsection}{\subsection}
\providecommand{\nobreakdash}{\nobreak\hbox{-}}
\setlength{\emergencystretch}{2em}
\multlinegap0pt
\usepackage{bm}
\usepackage{bbm}
\usepackage{dsfont}
\usepackage{xspace}
\usepackage{cancel}
\usepackage{mathtools}
\usepackage{amsthm}
\makeatletter
\@ifundefined{theorem}{\newtheorem{theorem}{Theorem}}{}
\@ifundefined{proposition}{\newtheorem{proposition}[theorem]{Proposition}}{}
\makeatother
\newcommand{\eS}{{\mathcal S}}
\newcommand{\hga}{{\hat\gamma}}
\title[Self-similar dendrites with finite boundary and P-sprouts]{Self-similar dendrites with finite boundary and P-sprouts}
\author{Andrei Tetenov, Ivan Yudin, Dmitrii Drozdov}
\date{Source: arXiv:2605.11833v3 (CC BY 4.0)}
\begin{document}
\maketitle

\noindent\emph{Source and licence.} Self-similar dendrites with finite boundary and P-sprouts. Version arXiv:2605.11833v3 (CC BY 4.0), distributed under the Creative Commons Attribution 4.0 International licence, as recorded in the arXiv licence metadata for this exact version and shown on the abstract page https://arxiv.org/abs/2605.11833v3. Full rights notice: licenses/arxiv-2605.11833-CC-BY-4.0.txt.\par\smallskip

\noindent\textit{Editorial scope.} Counted unit: the Proposition whose source label is \texttt{ord\_cut} in the article (source0.tex lines 476--484, source label \texttt{ord\_cut}), delivered together with the article's own complete original proof of it (source0.tex lines 486--514, one \texttt{proof} environment of 29 lines). No mathematics is written, repaired or re-derived: statement and proof are the article's own text line for line. Required context retained uncounted: none. Every definition, display and label of those lines is retained unchanged. Omitted from this package and not redistributed: 1--475, 485--485, 515--1277. Nothing inside a retained excerpt is omitted or altered: every retained body is delivered exactly as the article's lines are written, so no omission receipt of any kind accompanies this package. Citations: the retained excerpts carry no citation command at all, so no bibliography entry of the article is transcribed and no reference is redistributed. Reference adaptation: none was needed and none was made; every reference inside the delivered unit resolves inside it, so no unresolved reference or citation is produced. Printed slips kept literal: no printed word, symbol, hypothesis or number of the counted statement or of its proof was altered. Layout and class: the article's own class is \texttt{amsart} with packages including amsmath, amsthm, amsfonts, amssymb, graphicx, amscd, color, bm, hyperref, fontenc, inputenc, eucal; the wrapper is the lane's pinned \texttt{amsart} editorial wrapper at 10pt on A4, which restates the theorem-like environments the retained text needs and adds the article's own macro definitions that the retained text needs (\texttt{\textbackslash eS}, \texttt{\textbackslash hga}), with extra wrapper packages (bm, bbm, dsfont, xspace, cancel, mathtools). The delivered numbering is the unit's own. Assets: no retained span contains a figure environment, a graphics-inclusion command, a PDF-page inclusion, a table or a TikZ picture; no image file is copied into this package, the package declares no asset dependency and no image file is redistributed with it. Licence: version arXiv:2605.11833v3 of this article is distributed under the Creative Commons Attribution 4.0 International licence, as recorded in the arXiv licence metadata for this exact version and shown on the abstract page https://arxiv.org/abs/2605.11833v3. A full rights notice is delivered at licenses/arxiv-2605.11833-CC-BY-4.0.txt. Characters: every retained excerpt is pure ASCII, so the delivered engine, XeLaTeX with its default Latin Modern text font, reports no missing character.
\begin{proposition}\label{ord_cut}
Let the attractor $K$ of the system ${\eS}$ be a self-similar dendrite with finite boundary ${\partial} K$ and $\Gamma$ be its $P$-sprout.\\
1) %If $\Gamma({\eS})$ is regular, then 
For any $x\in K\setminus G_{\eS}(C)$, ${\mathop{\rm Ord} \nolimits}(x,K)\le\#P$.\\
2)  For any $x\in K\setminus G_{\eS}(C)$, there is ${\bf {i}}\in {I^*}$ such that ${\mathop{\rm Ord} \nolimits}(x,K)={\mathop{\rm Ord} \nolimits}(x,S_{\bf {i}}(\hga))$.\\
3) The set of cut points $CP(K)$ is contained in the set $ G_{\eS}(\hat \gamma)$.\\ 
4) Each  point $x\in K\setminus G_{{\eS}}(C)$ has a unique  address. If $x$ is a ramification point, this address is preperiodic.\\
5) $\hga{\subset} \bigcup\limits_{i=1}^m S_i(\hat \gamma)$.
\end{proposition}

\begin{proof} 1) Let  $x \in K\setminus G_{\eS}(C)$.  Let $Q_1,... Q_n$ be the connected components of $K \setminus \{x\}$. 
For each $k=1,...,n$ take  $x_k \in Q_k\setminus\{x\}$. There is ${\bf {i}}=i_1...i_n\in I^*$ such that
$x \in K_{\bf {i}}$ and
each $x_k\in K\setminus K_{\bf {i}}$. For each $k$ there is a unique arc $\gamma_k=\gamma(x,x_k){\subset} K$.


Note that $\gamma_k\setminus\{x\}{\subset} Q_k$ and
$\gamma_k\cap {\partial} K_{\bf {i}}=\{y_k\}$, where $y_k\neq x$. For different $k,l$,  
$\gamma(x,y_k)\cap \gamma(x,y_l)=\{x\}$, and all these arcs are  subarcs of
$S_{\bf {i}}(\hga)$. Therefore, $n\le\#P$.


To prove 2), take
$n={\mathop{\rm Ord} \nolimits}(x,K)$. Then ${\mathop{\rm Ord} \nolimits}(x,K)={\mathop{\rm Ord} \nolimits}(x,S_{\bf {i}}(\hga))$.

3) follows from 2). At the same time, $G_{\eS}(\hat \gamma){\supset} G_{\eS}({\partial} K)$; some points of the last set are the endpoints of $K$.

4) Since  $x\notin  G_{{\eS}}(C)$, the point $x$ has a unique address. 
By 2), there is a copy $K_{\bf {j}}$ such that  ${\mathop{\rm Ord} \nolimits}(x,K)={\mathop{\rm Ord} \nolimits}(x,S_{\bf {j}}(\hga))$. %For any boundary point $p\in{\partial} K$, the subarc $\gamma(x,p){\subset}\hga$ has a non-empty intersection 
%$\gamma(x,p)\cap{\partial} K_{\bf {j}}$, 
%therefore, 
%$ x\in \hga\cap K_{\bf {j}}{\subset} S_{\bf {j}}( \hga)$ and ${\mathop{\rm Ord} \nolimits}(S_{\bf {j}}^{-1}(x),\hga)\ge {\mathop{\rm Ord} \nolimits}(x,\hga)$. 
If $x$ is a ramification point,  its predecessor $ S_{\bf {j}}^{-1}(x)$ 
%$\in \hga\setminus G_{{\eS}}(C)$
is a ramification point of $\hga$ and lies in $\hga\setminus G_{{\eS}}(C)$.
Since the  number of  ramification points of $\hga$ is finite,  each ramification point $x\in K\setminus G_{{\eS}}({\partial} K)$ has a preperiodic address. 

5) Let $x\in\hga$ and $x\in S_i(K)$. If $x\notin {\partial} K_i$, there is a main subarc $\gamma(p_j,p_k)\ni x$. If any of these two points, say $p_j$, does not belong to ${\partial} K_i$, then $\gamma(x,p_j)\cap {\partial} K_i$ is a point  $\tilde p_j\in {\partial} K_i$. Thus, the arc $\gamma(p_j,p_k)\cap K_i$ is a  main subarc in  $K_i$. 
\end{proof}

\end{document}
